\documentclass[11pt]{article}

\usepackage[margin=1in]{geometry}
\usepackage{amsmath, amssymb}
\usepackage{algorithm}
\usepackage{algpseudocode}
\usepackage{hyperref}

\title{Artificial Neural Networks: Notes}
\author{ML A-Z --- Part 8}
\date{}

\begin{document}
\maketitle

\section{The Neuron}

An artificial neuron is a loose model of a biological neuron: it takes a
set of inputs, combines them, and fires an output signal. For an input
vector $\mathbf{x} = (x_1, \dots, x_m)$ (e.g. the independent variables of
one observation, typically standardized so all inputs are on comparable
scales), a neuron computes

\[
  z = \sum_{i=1}^{m} w_i x_i + b = \mathbf{w}^\top \mathbf{x} + b,
\]

where $w_i$ is the \textbf{weight} on input $x_i$ and $b$ is the
\textbf{bias} (an always-on input, sometimes written as $w_0 x_0$ with
$x_0 = 1$). The weighted sum $z$ is then passed through an
\textbf{activation function} $\phi$ to produce the neuron's output
$a = \phi(z)$.

Weights are the parameters the network learns: they encode how strongly
each input signal influences the neuron downstream.

\section{Activation Functions}

Without an activation function, a network is just a composition of linear
maps, which collapses to a single linear map --- no matter how many layers
are stacked, the network could only ever learn linear relationships.
Activation functions introduce nonlinearity, which is what lets a network
approximate arbitrary functions.

\begin{itemize}
  \item \textbf{Threshold function:}
        \[
          \phi(z) = \begin{cases} 1 & z \ge 0 \\ 0 & z < 0 \end{cases}
        \]
        A hard on/off gate. Not differentiable at $0$, so it cannot be
        used with gradient-based learning.

  \item \textbf{Sigmoid:}
        \[
          \phi(z) = \frac{1}{1 + e^{-z}}
        \]
        Smooth, squashes $z$ into $(0, 1)$. Useful in the output layer for
        binary classification, since the output can be read as a
        probability. Suffers from \emph{vanishing gradients}: for large
        $|z|$, $\phi'(z) \approx 0$, so learning slows down in saturated
        neurons.

  \item \textbf{Rectifier (ReLU):}
        \[
          \phi(z) = \max(0, z)
        \]
        The default choice for hidden layers in modern networks. Cheap to
        compute, and its gradient is either $0$ or $1$ (no saturation on
        the positive side), which mitigates vanishing gradients.

  \item \textbf{Hyperbolic tangent (tanh):}
        \[
          \phi(z) = \frac{e^z - e^{-z}}{e^z + e^{-z}} \in (-1, 1)
        \]
        Like sigmoid but zero-centered, which tends to make optimization
        better behaved.
\end{itemize}

A common rule of thumb: ReLU in hidden layers, sigmoid (binary) or softmax
(multi-class) at the output layer.

\section{Network Architecture: Forward Propagation}

A fully-connected feed-forward network arranges neurons in layers: an
input layer, one or more \textbf{hidden layers}, and an output layer.
Every neuron in layer $\ell$ receives, as its inputs, the outputs of every
neuron in layer $\ell - 1$.

For layer $\ell$ with weight matrix $W^{(\ell)}$ and bias vector
$\mathbf{b}^{(\ell)}$,

\[
  \mathbf{z}^{(\ell)} = W^{(\ell)} \mathbf{a}^{(\ell-1)} + \mathbf{b}^{(\ell)},
  \qquad
  \mathbf{a}^{(\ell)} = \phi\!\left(\mathbf{z}^{(\ell)}\right),
\]

with $\mathbf{a}^{(0)} = \mathbf{x}$ (the raw input). Applying this
layer-by-layer from input to output is \textbf{forward propagation}: it
turns an input observation into a prediction $\hat{y} = \mathbf{a}^{(L)}$
for a network of $L$ layers.

Intuitively, each hidden layer re-represents the data in a new feature
space; deeper layers can combine the simple patterns detected by earlier
layers into more complex, abstract features.

\section{Cost Function}

Learning means choosing weights and biases so that the network's
predictions $\hat{y}$ are close to the true targets $y$. This is measured
with a \textbf{cost function} (loss function). A standard choice for
regression is squared error,

\[
  C = \frac{1}{2}\left(\hat{y} - y\right)^2,
\]

summed or averaged over all training examples. (For classification,
cross-entropy is typically used instead, since it pairs naturally with a
sigmoid/softmax output and produces better gradients when predictions are
confidently wrong.)

Training is the optimization problem of finding the weights $\mathbf{w}$
that minimize $C(\mathbf{w})$ over the training set.

\section{Gradient Descent}

Since $C$ is generally a nonconvex, high-dimensional function of every
weight and bias in the network, it cannot be minimized in closed form.
\textbf{Gradient descent} minimizes it iteratively: at each step, nudge
every parameter in the direction that decreases $C$ fastest, i.e. opposite
its gradient,

\[
  \mathbf{w} \gets \mathbf{w} - \eta \, \nabla_{\mathbf{w}} C,
\]

where $\eta$ is the \textbf{learning rate}. Too large an $\eta$ can cause
the update to overshoot and diverge; too small makes training
impractically slow.

\subsection{Batch vs. Stochastic Gradient Descent}

\begin{itemize}
  \item \textbf{Batch gradient descent} computes $\nabla_{\mathbf{w}} C$
        using the \emph{entire} training set before making one update.
        This gives an accurate gradient, but each step is expensive, and
        because $C$ is generally nonconvex, batch GD can get stuck in a
        local minimum.

  \item \textbf{Stochastic gradient descent (SGD)} updates the weights
        after \emph{each} individual training example, using that
        example's (noisy) gradient. The noise in these single-example
        updates actually helps: it lets the optimizer jump out of local
        minima that would trap batch GD, and it makes each step far
        cheaper. SGD is also better suited to large datasets, since it
        does not need to hold the whole dataset in memory to take a step.

  \item \textbf{Mini-batch gradient descent} is the practical middle
        ground used by most deep learning frameworks: compute the
        gradient over a small batch (e.g. 32--256 examples) at a time.
        This retains much of SGD's escape-local-minima behavior and
        memory efficiency, while vectorizing the batch's computation for
        speed.
\end{itemize}

\section{Backpropagation}

Backpropagation is the algorithm that computes $\nabla_{\mathbf{w}} C$
efficiently, layer by layer, using the chain rule --- rather than
perturbing each weight independently (which would cost one full forward
pass per weight).

\subsection{Setup}

Let layer $\ell$ have weighted input $\mathbf{z}^{(\ell)}$ and activation
$\mathbf{a}^{(\ell)} = \phi(\mathbf{z}^{(\ell)})$, for $\ell = 1, \dots,
L$, with $\mathbf{a}^{(L)} = \hat{y}$ the network output. Define the
\textbf{error} at layer $\ell$ as

\[
  \boldsymbol{\delta}^{(\ell)} = \frac{\partial C}{\partial \mathbf{z}^{(\ell)}}.
\]

\subsection{Output layer error}

By the chain rule,

\[
  \boldsymbol{\delta}^{(L)}
  = \frac{\partial C}{\partial \mathbf{a}^{(L)}} \odot \phi'\!\left(\mathbf{z}^{(L)}\right),
\]

where $\odot$ is the elementwise product. For squared error,
$\partial C / \partial \mathbf{a}^{(L)} = \hat{y} - y$.

\subsection{Propagating the error backward}

The error at layer $\ell$ depends on the error one layer ahead of it:

\[
  \boldsymbol{\delta}^{(\ell)}
  = \left(\left(W^{(\ell+1)}\right)^{\top} \boldsymbol{\delta}^{(\ell+1)}\right)
    \odot \phi'\!\left(\mathbf{z}^{(\ell)}\right).
\]

This is the key equation: it lets $\boldsymbol{\delta}^{(\ell)}$ be computed
from $\boldsymbol{\delta}^{(\ell+1)}$, moving backward from the output
layer to the input layer --- hence \emph{back}propagation. It is exactly
the chain rule applied once per layer, reusing work already done for the
layer ahead instead of recomputing it from scratch.

\subsection{Gradients from the error}

Once every $\boldsymbol{\delta}^{(\ell)}$ is known, the gradients needed
for the weight update fall out directly:

\[
  \frac{\partial C}{\partial W^{(\ell)}} = \boldsymbol{\delta}^{(\ell)} \left(\mathbf{a}^{(\ell-1)}\right)^{\top},
  \qquad
  \frac{\partial C}{\partial \mathbf{b}^{(\ell)}} = \boldsymbol{\delta}^{(\ell)}.
\]

\subsection{Algorithm}

\begin{algorithm}[h]
\caption{Training one example with backpropagation + SGD}
\begin{algorithmic}[1]
\State \textbf{Forward pass:} compute $\mathbf{z}^{(\ell)}, \mathbf{a}^{(\ell)}$ for $\ell = 1, \dots, L$
\State \textbf{Output error:} $\boldsymbol{\delta}^{(L)} \gets (\hat{y} - y) \odot \phi'(\mathbf{z}^{(L)})$
\For{$\ell = L-1, \dots, 1$}
  \State $\boldsymbol{\delta}^{(\ell)} \gets \left((W^{(\ell+1)})^\top \boldsymbol{\delta}^{(\ell+1)}\right) \odot \phi'(\mathbf{z}^{(\ell)})$
\EndFor
\For{$\ell = 1, \dots, L$}
  \State $W^{(\ell)} \gets W^{(\ell)} - \eta\, \boldsymbol{\delta}^{(\ell)} (\mathbf{a}^{(\ell-1)})^\top$
  \State $\mathbf{b}^{(\ell)} \gets \mathbf{b}^{(\ell)} - \eta\, \boldsymbol{\delta}^{(\ell)}$
\EndFor
\end{algorithmic}
\end{algorithm}

Repeating this over all training examples, for many passes (\textbf{epochs}),
is how the network learns: each epoch, weights move a little further in
the direction that reduces the cost function.

\section{Key Takeaways}

\begin{itemize}
  \item Activation functions are what give a network the capacity to
        learn nonlinear relationships; ReLU is the standard hidden-layer
        default, sigmoid/softmax the standard output-layer choice.
  \item Forward propagation produces a prediction; backpropagation
        produces the gradient of the cost with respect to every weight,
        by reusing per-layer error terms via the chain rule instead of
        recomputing each partial derivative independently.
  \item SGD (or mini-batch GD) trades a noisier gradient estimate for
        much cheaper, much more frequent updates, and that noise is what
        helps escape local minima in the nonconvex cost surface.
  \item The learning rate $\eta$ controls the step size and is the main
        knob balancing convergence speed against stability.
\end{itemize}

\end{document}
